% 第19批：学习指导自测题8选择题8.1.1--8.1.8；印刷171--172／PDF179--180。
\subsection{自测题8：电容、磁场叠加和位移电流选择题}

\sourcetest{8.1.1}{171}{179}
\problemtext 三只大小相同的金属小球，球1和2各带等量同号电荷 $q$，相距远于球径，库仑斥力为 $F$。中性球3先与1接触、再与2接触后移走，求球1与球2间的斥力。
\workings 首次接触后，球1和3各带 $q/2$。球3再与球2接触时总电量为 $3q/2$，两球各得 $3q/4$。最终
\[\boxed{F'=\frac{k(q/2)(3q/4)}{r^2}=\frac38F}.\]
选 \textbf{D}。
\crosscheck 第二次接触时球3的电荷不是零，而是第一次接触后的 $q/2$；不能误作 $F/4$。

\sourcetest{8.1.2}{171}{179}
\problemtext 半径为 $R_1,R_2$（$R_2>R_1$）的两同心金属薄球壳最初分别带电 $q_1,q_2$，外壳初始电势记 $U_2$（无穷远为零）。接导线使两壳等势后，求共同电势。
\workings 两壳导通后构成一个整体导体，无电荷保留在无内部电荷的腔内表面，所有净电荷 $q_1+q_2$ 分布在半径 $R_2$ 的最外表面。因此
\[U_{\rm common}=\frac{q_1+q_2}{4\pi\varepsilon_0R_2}.\]
而外壳导通前已满足 $U_2=(q_1+q_2)/(4\pi\varepsilon_0R_2)$，故
\[\boxed{U_{\rm common}=U_2}.\]
选 \textbf{B}。
\crosscheck 内壳初始电势 $U_1$ 不一定等于 $U_2$；导线相连使电荷流动，不能把两壳视为始终保持原始各自电荷。

\sourcetest{8.1.3}{171}{179}
\problemtext 电容器 $C_1=8\,\mu{\rm F}$、$C_2=2\,\mu{\rm F}$ 分别充至1000V，然后按题8.1.1图反极性并联，求两电容器最终两接点的电势差大小。
\workings 并联共用的正向节点原净自由电荷为 $Q_0=C_1U_0-C_2U_0$，总电容量为 $C_1+C_2$。电荷守恒得
\[
\boxed{U=\frac{C_1-C_2}{C_1+C_2}U_0
=\frac{8-2}{8+2}(1000)=600\,{\rm V}}.
\]
选 \textbf{C}。
\crosscheck 两电容器接线的原极性相反，电荷应相减而非相加；实际再分配过程中的能量损失不影响最终电荷平衡式。

\sourcetest{8.1.4}{171}{179}
\problemtext 题8.1.2图中电流从无穷长竖直引线1在顶点 $a$ 流入由四条等阻边组成的正方形框，再从对角顶点 $c$ 经引线2流出；比较框中心 $O$ 处引线场与方框场。
\workings 两条从 $a$ 到 $c$ 的边路径电阻相同，电流各为 $I/2$。正方形四边中相对的两支路对 $O$ 的场大小相等、方向相反，故 $\boxed{B_3=0}$。两条无限长共线延伸引线1、2在 $O$ 产生大小相同方向相反的场，故 $\boxed{\bm B_1+\bm B_2=0}$，但 $\bm B_1,\bm B_2$ 各不为零。总场
\[\boxed{\bm B_O=\bm0}.\]
选 \textbf{B}。
\crosscheck 总磁场为零不意味着每根引线贡献都为零；须把相互抵消与单独为零区分。

\sourcetest{8.1.5}{171}{179}
\problemtext 在 $x=1$ 和 $x=3$ 的两条无限长平行导线均沿 $+y$ 方向载同一电流 $I$。问磁场大小为零的点位于何处。
\workings 两根直导线在 $x$--$y$ 平面内、位于两线之间的场方向相反；其大小为
\[B_1=\frac{\mu_0I}{2\pi|x-1|},\quad
B_2=\frac{\mu_0I}{2\pi|x-3|}.\]
相等须 $|x-1|=|x-3|$，故
\[\boxed{x=2}\]
的平面内中线处磁场相消，选 \textbf{A}。
\crosscheck 严格地说只有位于两根导线共同张成的平面、且两线等距的那条直线 $x=2,z=0$ 上场为零；不能说整个三维平面 $x=2$ 的所有点均为零。

\sourcetest{8.1.6}{172}{180}
\problemtext 题8.1.4图：直角三角形金属框 $abc$ 绕竖直边 $ab$ 转动，边 $bc=\ell$ 且 $ab\parallel\bm B$。求整闭合回路动生电动势与 $U_a-U_c$。
\workings 由于转轴与均匀磁场平行，线框转动时始终有 $\bm B$ 位于线框平面中，磁通量恒为零，故
\[\boxed{\mathcal E_{\rm loop}=0}.\]
若角速度矢量与 $\bm B$ 同向，则正电荷受 $\bm v\times\bm B$ 的径向力，距轴 $\ell$ 的 $c$ 电势高于轴上 $a$：
\[\boxed{U_a-U_c=-\int_0^\ell B\omega r\,dr=-\tfrac12B\omega\ell^2}.\]
选 \textbf{B}（按原图角速度方向）。
\crosscheck 闭合回路电动势为零不代表两端点无电势差；若角速度方向反向，端点电势差符号也反向。

\sourcetest{8.1.7}{172}{180}
\problemtext 相距 $2a$ 的两无限长平行直导线组成闭合回路，中间电流方向相反，电流大小均为 $I$。求中点 $P$ 的磁能密度。
\workings 两导线在中点产生的磁场方向相同，故
\[B_P=2\frac{\mu_0I}{2\pi a}=\frac{\mu_0I}{\pi a}.\]
真空磁场能量密度为
\[\boxed{w_B=\frac{B_P^2}{2\mu_0}
=\frac1{2\mu_0}\left(\frac{\mu_0I}{\pi a}\right)^2}.\]
选 \textbf{C}。
\crosscheck 电流沿两线方向相反，线间的磁场叠加而非抵消；不可把能量密度先算两个单独项再相加以忽略交叉项。

\sourcetest{8.1.8}{172}{180}
\problemtext 空气平行板电容器两板电势差随时间按 $U_{ab}=Kt$（$K>0$）增长。两板间有点1和点2，其中2更靠近板的外缘，比较两点感生磁场大小 $B_1,B_2$。
\workings 忽略边缘效应，板间 $E=Kt/d$、位移电流密度 $J_d=\varepsilon_0K/d$ 均匀。对距极板中心轴 $r$ 的圆形安培环路，
\[B_\varphi(r)2\pi r=\mu_0J_d\pi r^2
\Longrightarrow\boxed{B_\varphi(r)=\frac{\mu_0\varepsilon_0 K}{2d}\,r}.\]
故 $r_2>r_1$ 时
\[\boxed{B_1<B_2},\]
选 \textbf{B}。
\crosscheck 比较的是随时间变化的电场产生的磁场，极板间无须真实导电流穿过介质；近板缘时边缘场会修正该线性关系。
